% Part: incompleteness
% Chapter: arithmetization-syntax
% Section: coding-formulas

\documentclass[../../../include/open-logic-section]{subfiles}

\begin{document}

\olfileid{inc}{art}{frm}
\olsection{\printtoken{P}{formula}નું સંકેતબદ્ધીકરણ}

સંબંધ~$\fn{Term}(x)$ને આદિમ
પુનરાવર્તી રીતે વ્યાખ્યાયિત કર્યા પછી,
તેની મદદથી !!{formula}s માટેનો અનુરૂપ
સંબંધ~$\fn{Frm}(x)$ પણ આદિમ
પુનરાવર્તી રીતે વ્યાખ્યાયિત કરી શકીએ.

\begin{prop}
$x$ આણ્વિક !!{formula}ની ગડેલ-સંખ્યા
હોય ત્યારે અને માત્ર ત્યારે સત્ય થતો
સંબંધ~$\fn{Atom}(x)$ આદિમ પુનરાવર્તી છે.
\end{prop}

\begin{proof}
નીચેમાંથી એક શરત હોય ત્યારે અને માત્ર
ત્યારે નંબર~$x$ આણ્વિક !!{formula}ની
ગડેલ-સંખ્યા છે:
\begin{enumerate}
\item એવાં~$n$, $j < x$ અને~$z < x$
  હોય કે દરેક~$i < n$ માટે
  $\fn{Term}((z)_i)$ અને $x = $
\[
\Gn{\Obj P^n_j(} \concat \fn{flatten}(z) \concat \Gn{)}.
\]
\item એવાં~$z_1, z_2 < x$ હોય કે
  $\fn{Term}(z_1)$, $\fn{Term}(z_2)$
  અને $x = $
\[
\Gn{{\eq}(} \concat z_1 \concat \Gn{,} \concat z_2 \concat{\Gn{)}}.
\]
  \tagitem{prvFalse}{$x = \Gn{\lfalse}$.}{}
  \tagitem{prvTrue}{$x = \Gn{\ltrue}$.}{}
\end{enumerate}
\end{proof}

\begin{prop}
\ollabel{prop:frm-primrec}
$x$ કોઈ !!a{formula}ની ગડેલ-સંખ્યા
હોય ત્યારે અને માત્ર ત્યારે સત્ય થતો
સંબંધ~$\fn{Frm}(x)$ આદિમ પુનરાવર્તી છે.
\end{prop}

\begin{proof}
ચિહ્નોનો અનુક્રમ~$s$ !!a{formula} છે
ત્યારે અને માત્ર ત્યારે, જ્યારે
!!{formula}ની એવી રચના-શ્રેણી~$s_0$,
\dots, $s_{k-1} = s$ હોય જે રચના-નિયમો
મુજબ આણ્વિક !!{formula}sમાંથી~$s$
કેવી રીતે બન્યું તે નોંધે છે. બિનજરૂરી
વચગાળાનાં સૂત્રો કાઢીને માત્ર ઉપસૂત્રોની
રચના-શ્રેણી પસંદ કરીએ, તો દરેક~$s_i$નો
સંકેત નંબર અંતિમ~$s$ના સંકેત નંબર~$x$થી
મોટો નથી. સમગ્ર રચના-શ્રેણી
$\tuple{s_0, \dots, s_{k-1}}$ના સંકેત
નંબર માટે પણ અંતિમ નંબરથી ગણાતી
આદિમ પુનરાવર્તી મર્યાદા મળે છે.
\sourcecorrection{OLINC-005}{સ્થિર અંગ્રેજી મૂળમાં
અંતિમ સભ્ય સહિત દરેક~$s_i$નો સંકેત
નંબર અને આખી રચના-શ્રેણીનો સંકેત
નંબર બંને~$x$ કરતાં નાના કહેવાયા છે.
અંતિમ સભ્યનો સંકેત પોતે જ~$x$ છે,
અને આખી શ્રેણીનો સંકેત તેનાથી મોટો
થઈ શકે છે. બિનજરૂરી વચગાળાનાં સૂત્રો
કાઢીને ઉપસૂત્રો પૂરતી રચના-શ્રેણી લેવી
પડે. પહેલાંના પદ-પ્રમાણની જેમ
અહીં જરૂરી વાત~$x$ પરથી ગણાતી આદિમ
પુનરાવર્તી મર્યાદા હોવાની છે.}
\end{proof}

\begin{prob}
પહેલી સાબિતીની રીત અનુસરીને
\olref[inc][art][frm]{prop:frm-primrec}ની
વિગતવાર સાબિતી આપો;
\olref[inc][art][trm]{prop:term-primrec}ની
પહેલી સાબિતી માર્ગદર્શક છે.
\end{prob}

\begin{prop}
\ollabel{prop:freeocc-primrec}
ગડેલ-સંખ્યા~$x$ ધરાવતા સૂત્રનું
$i$મું ચિહ્ન એ ગડેલ-સંખ્યા~$z$
ધરાવતા ચલની મુક્ત ઘટના હોય ત્યારે
અને માત્ર ત્યારે સત્ય થતો સંબંધ
$\fn{FreeOcc}(x, z, i)$ આદિમ
પુનરાવર્તી છે.
\end{prop}

\begin{proof}
અભ્યાસપ્રશ્ન.
\end{proof}

\begin{prob}
\olref[inc][art][frm]{prop:freeocc-primrec}
સાબિત કરો. !!a{formula}નો કોઈ
ઉપઅનુક્રમ પણ !!a{formula} હોય, તો
તે મૂળ સૂત્રનું ઉપ-!!{formula} હોય છે,
એ તથ્ય વાપરી શકો છો.
\end{prob}

\begin{prop}
$x$ કોઈ !!a{sentence}ની ગડેલ-સંખ્યા
હોય ત્યારે અને માત્ર ત્યારે સત્ય થતો
ગુણધર્મ~$\fn{Sent}(x)$ આદિમ
પુનરાવર્તી છે.
\end{prop}

\begin{proof}
!!{sentence} એ !!{variable}sની મુક્ત
ઘટના વિનાનું !!a{formula} છે. તેથી
$\fn{Sent}(x)$ ત્યારે અને માત્ર ત્યારે,
જ્યારે
\begin{multline*}
\bforall{i<\len{x}}{\bforall{z<x}{}}\\
(\bexists{j<z}{z=\Gn{\Obj v_j}} \lif \lnot\fn{FreeOcc}(x,z,i)).
\end{multline*}
\end{proof}


\end{document}
